% ============================================================
% Contrôle : Équations différentielles du premier ordre
% BTS MS 2e année - Sujet D (mardi 3 novembre 2026, 15h15)
% Pôle Formation UIMM - CVDL - S. JAUBERT
% Corrigé enseignant : MOD06_EDL1_Controle_MS_Corrige_v1_0
% Compilation : pdflatex, deux passes (nombre total de pages).
% ============================================================
\documentclass[a4paper,11pt]{article}
\usepackage[utf8]{inputenc}
\usepackage[T1]{fontenc}
% babel-french absent sur certains postes : chargement conditionnel
\IfFileExists{french.ldf}{\usepackage[french]{babel}}{}
\usepackage{helvet}
\renewcommand{\familydefault}{\sfdefault}
\usepackage[top=1.8cm, bottom=3cm, left=2cm, right=2cm, headheight=14pt, footskip=1.6cm]{geometry}
\usepackage{amsmath,amssymb}
\usepackage{graphicx}
\usepackage[dvipsnames,table]{xcolor}
\usepackage[most]{tcolorbox}
\usepackage{fancyhdr}
\usepackage{enumitem}
\usepackage{array,tabularx,needspace}
\usepackage{lastpage}
\usepackage{hyperref}

% ----- Charte UIMM -----
\definecolor{uimmblue}{HTML}{005C85}
\definecolor{uimmred}{HTML}{D31326}
\definecolor{uimmlightblue}{HTML}{E8F1F8}
\definecolor{uimmgray}{HTML}{333333}
\definecolor{uimmlightgray}{HTML}{F5F5F5}
\hypersetup{colorlinks=true, linkcolor=uimmblue, urlcolor=uimmblue}
\setlength{\parindent}{0pt}
\setlength{\parskip}{3pt}
\setlist[enumerate,1]{label={\color{uimmblue}\textbf{\arabic*.}}, leftmargin=*, itemsep=3pt}
\setlist[enumerate,2]{label=\alph*), itemsep=2pt}

% ----- Paramètres du sujet -----
\newcommand{\groupe}{D}
\newcommand{\creneau}{Mardi 3 novembre 2026, 15h15}

% ----- En-tête et pied de page (modèle PR3) -----
\pagestyle{fancy}
\fancyhf{}
\renewcommand{\headrulewidth}{0pt}
\fancyfoot[C]{\small\renewcommand{\arraystretch}{1.3}%
\begin{tabularx}{\textwidth}{|>{\centering\arraybackslash\itshape}X|>{\centering\arraybackslash}p{1.8cm}|>{\centering\arraybackslash}p{3.6cm}|>{\centering\arraybackslash}p{1.9cm}|}
\hline
BTS MS 2 -- MOD06 Équations différentielles -- Contrôle, sujet \groupe & {\scriptsize\itshape Auteur :}\newline SJ & MAJ du 01/11/2026\newline{\color{uimmblue}Version 1.0} & Page :\newline \thepage/\pageref*{LastPage}\\
\hline
\end{tabularx}}

% ----- Exercices -----
\newcounter{exo}
% \exo{titre}{compétences}{capacités}{barème}
\newcommand{\exo}[4]{\needspace{0.5\textheight}\refstepcounter{exo}\vspace{8pt}%
\begin{tcolorbox}[colback=uimmlightblue, colframe=uimmblue, boxrule=0.8pt, arc=2pt, left=6pt, right=6pt, top=3pt, bottom=3pt]
{\color{uimmblue}\bfseries Exercice \theexo\ : #1}\hfill{\color{uimmblue}\bfseries #4}\\
{\small\color{uimmred}\bfseries #2\quad #3}
\end{tcolorbox}}
% Espace de réponse (blanc), hauteur en cm
\newcommand{\rep}[1]{\par\noindent\begin{minipage}[t][#1cm]{\linewidth}\mbox{}\end{minipage}\par}
\newcommand{\pts}[1]{\unskip\nobreak\hfill\mbox{\small\color{uimmgray}(#1)}}

\begin{document}
\thispagestyle{fancy}

% ============================================================
% EN-TÊTE PR3
% ============================================================
\noindent{\renewcommand{\arraystretch}{1.1}%
\begin{tabularx}{\textwidth}{|>{\centering\arraybackslash}m{2.9cm}|>{\centering\arraybackslash}m{4.3cm}|>{\centering\arraybackslash}X|}
\hline
\rule{0pt}{2.4cm}\IfFileExists{logo_uimm_placeholder.jpg}{\includegraphics[height=2.2cm]{logo_uimm_placeholder.jpg}}{\fbox{\parbox[c][2cm][c]{2.4cm}{\centering\bfseries\color{uimmblue}UIMM}}} &
{\bfseries\large CFAI CENTRE}\newline{\footnotesize 74, Route Nationale\newline La Chapelle St Mesmin\newline orleans@poleformation-\newline uimmcvdl.fr} &
{\Large\bfseries\color{uimmblue}Contrôle : Équations différentielles\newline{}du premier ordre}\newline{}{\small BTS MS 2\textsuperscript{e} année -- \textbf{Sujet \groupe}}\newline{}{\footnotesize\creneau}\\
\hline
\end{tabularx}}

\vspace{-1pt}
\noindent{\footnotesize\renewcommand{\arraystretch}{1.2}\hyphenpenalty=10000\exhyphenpenalty=10000%
\newcommand{\on}{\cellcolor{uimmblue!22}}%
\arrayrulecolor{black}\setlength{\arrayrulewidth}{0.7pt}%
\newcolumntype{C}[1]{>{\centering\arraybackslash}m{#1}}%
\begin{tabular}{|C{1.85cm}|C{2.05cm}|C{1.85cm}|C{2.0cm}|C{1.8cm}|C{2.15cm}|C{2.2cm}|}
\hline
\bfseries C1 S'informer & \bfseries C2 Chercher & \on\bfseries C3 Modéliser & \on\bfseries C4 Raisonner, argumenter & \multicolumn{2}{C{4.38cm}|}{\cellcolor{uimmblue!22}\bfseries C5 Calculer, illustrer, mettre en œuvre une démarche} & \on\bfseries C6 Communiquer\\
\hline
\end{tabular}}

\vspace{8pt}
% ----- Cartouche élève -----
\noindent{\renewcommand{\arraystretch}{2.1}%
\begin{tabularx}{\textwidth}{|>{\bfseries\color{uimmblue}}l X|>{\bfseries\color{uimmblue}}l p{2.6cm}|>{\bfseries\color{uimmblue}}l p{1.6cm}|}
\hline
Nom : & & Groupe : & BTS MS 2 -- \groupe & Note : & \hfill / 20\\
\hline
Prénom : & & \multicolumn{4}{l|}{\color{uimmgray}\footnotesize Répondre dans les espaces prévus.}\\
\hline
\end{tabularx}}

\vspace{6pt}
\begin{tcolorbox}[colback=uimmlightgray, colframe=uimmgray, boxrule=0.6pt, arc=2pt, left=8pt, right=8pt, top=4pt, bottom=4pt]
\small\textbf{Durée : 1 h 15.} Aucun document autorisé. Téléphone portable interdit (éteint, rangé dans le sac). Calculatrice autorisée en \textbf{mode examen}. Répondre \textbf{directement sur le sujet}, dans les espaces prévus. Toute réponse doit être justifiée par un calcul. Barème indicatif sur 20.
\end{tcolorbox}

% ============================================================
% RAPPEL
% ============================================================
\begin{tcolorbox}[colback=uimmlightblue!40, colframe=uimmblue, boxrule=0.8pt, arc=2pt, left=8pt, right=8pt, top=4pt, bottom=4pt, title={\bfseries Rappel : résoudre $a\,y' + b\,y = c$ ($a$, $b$, $c$ constantes, $a \neq 0$)}, coltitle=white, colbacktitle=uimmblue]
\small
\begin{enumerate}[leftmargin=*, itemsep=1pt, topsep=1pt]
  \item \textbf{Réduire l'équation} : tous les termes en $y'$ et en $y$ à gauche, la constante à droite. Lire $a$, $b$ et $c$.
  \item \textbf{Forme normalisée} : diviser par $a$, on obtient $y' + \alpha\,y = \dfrac{c}{a}$ avec $\alpha = \dfrac{b}{a}$.
  \item \textbf{Équation homogène} $y' + \alpha\,y = 0$ : ses solutions sont $y_h(t) = C\,e^{-\alpha t}$, $C$ constante réelle.
  \item \textbf{Solution particulière constante} (si $b \neq 0$) : $y_p = \dfrac{c}{b}$.
  \item \textbf{Solution générale} : $y(t) = y_h(t) + y_p = C\,e^{-\alpha t} + \dfrac{c}{b}$.
  \item \textbf{Condition initiale} $y(t_0) = y_0$ : remplacer $t$ par $t_0$, en déduire $C$, écrire la solution unique.
\end{enumerate}
\textbf{Constante de temps} (si $\alpha > 0$) : $\tau = \dfrac{1}{\alpha}$. À $t = \tau$, environ $63\,\%$ de l'évolution est réalisée.\\
\textbf{Exponentielle et logarithme} : $e^{0} = 1$ ; \ $\big(e^{kt}\big)' = k\,e^{kt}$ ; \ si $\alpha > 0$, $e^{-\alpha t} \to 0$ quand $t \to +\infty$ ; \ pour $A > 0$ : $e^{X} = A \iff X = \ln A$ ; \ $\ln\dfrac{1}{A} = -\ln A$.
\end{tcolorbox}

% ============================================================
\exo{Résolution d'équations différentielles}{C4 C5}{[S01] [C01] [C02]}{11 points}
% ============================================================
La fonction inconnue $y$ est une fonction de la variable $t$. Pour chaque équation : la réduire, la résoudre, puis déterminer la solution qui vérifie la condition initiale.

\medskip
\textbf{\color{uimmblue}Partie A : équations homogènes}
\begin{enumerate}
  \needspace{5.1cm}
  \item $2y' = 2y$ \quad et \quad $y(0) = -2$. \pts{1 pt}
  \rep{3.5}
  \needspace{5.6cm}
  \item $4y' + 3y = y' - 3y$ \quad et \quad $y(0) = 3$. \pts{1,5 pt}
  \rep{4.0}
  \needspace{5.1cm}
  \item $y = 2y'$ \quad et \quad $y(0) = 5$. \pts{1 pt}
  \rep{3.5}
  \needspace{6.1cm}
  \item $\dfrac{y'}{3} + y = 0$ \quad et \quad $y(1) = 4$ \quad (la condition n'est pas donnée en $t = 0$). \pts{1,5 pt}
  \rep{4.5}
\end{enumerate}

\needspace{7cm}\medskip
\textbf{\color{uimmblue}Partie B : second membre constant}
\begin{enumerate}[resume]
  \needspace{5.6cm}
  \item $y' + 3y = 12$ \quad et \quad $y(0) = 0$. \pts{1 pt}
  \rep{4.0}
  \needspace{6.1cm}
  \item $4y' + 2y = 2y' - 2y + 12$ \quad et \quad $y(0) = 1$. \pts{1,5 pt}
  \rep{4.5}
  \needspace{6.1cm}
  \item $\dfrac{y'}{2} = \dfrac{y'}{5} + 3$ \quad et \quad $y(0) = 2$. \pts{1,5 pt}
  \rep{4.5}
  \needspace{6.6cm}
  \item $4\,(y - 3) = 2y'$ \quad et \quad $y(0) = 5$. \pts{2 pts}
  \rep{5.0}
\end{enumerate}

% ============================================================
\exo{Passage en vitesse nominale d'un convoyeur à bande}{C3 C4 C5 C6}{[S01] [C01] [C02]}{9 points}
% ============================================================
Sur une ligne de conditionnement, un convoyeur à bande tourne en vitesse d'approche $0{,}3$\,m/s. À l'instant $t = 0$, l'automate commande le passage en vitesse nominale $V_N = 1{,}5$\,m/s. On note $v(t)$ la vitesse de la bande (en m/s), $t$ étant le temps en secondes. Ainsi $v(0) = 0{,}3$.

On admet le modèle suivant : l'accélération $v'(t)$ de la bande est \textbf{proportionnelle} à l'écart entre la vitesse nominale $V_N$ et la vitesse $v(t)$, avec le coefficient de proportionnalité $k = 0{,}8\ \text{s}^{-1}$.

Les cartons ne peuvent être déposés sur la bande que si sa vitesse est supérieure ou égale à $1{,}4$\,m/s (seuil retenu pour cet exercice).

\begin{enumerate}
  \needspace{4.6cm}
  \item Traduire le modèle par une égalité entre $v'(t)$ et $v(t)$. Montrer que $v$ vérifie l'équation \[ (E) : \quad v' + 0{,}8\,v = 1{,}2 . \] \pts{1 pt}
  \rep{3.0}
  \needspace{3.1cm}
  \item Résoudre l'équation homogène associée à $(E)$. \pts{0,5 pt}
  \rep{1.5}
  \needspace{3.1cm}
  \item Déterminer une solution particulière constante de $(E)$. \pts{0,5 pt}
  \rep{1.5}
  \needspace{5.6cm}
  \item En déduire la solution générale de $(E)$, puis montrer que $v(t) = 1{,}5 - 1{,}2\,e^{-0{,}8\,t}$. \pts{1,5 pt}
  \rep{4.0}
  \needspace{4.6cm}
  \item Calculer $v'(t)$. Étudier son signe et en déduire le sens de variation de $v$ sur $[0\,;\,+\infty[$. Interpréter. \pts{1 pt}
  \rep{3.0}
  \needspace{3.6cm}
  \item Déterminer la limite de $v(t)$ quand $t$ tend vers $+\infty$. Interpréter ce résultat pour le convoyeur. \pts{1 pt}
  \rep{2.0}
  \needspace{4.6cm}
  \item On note $\tau = \dfrac{1}{0{,}8}$ la constante de temps. Calculer $\tau$, puis $v(\tau)$ arrondi au centième. Quel pourcentage de l'accélération totale (de $0{,}3$\,m/s à la limite) est réalisé à cet instant ? \pts{1 pt}
  \rep{3.0}
  \needspace{5.6cm}
  \item Déterminer par le calcul l'instant $t_1$ à partir duquel les cartons peuvent être déposés. Arrondir au centième de seconde. \pts{1,5 pt}
  \rep{4.0}
  \needspace{4.6cm}
  \item L'automate est programmé pour autoriser le dépôt des cartons $3$ secondes après la commande. Ce réglage est-il correct ? Justifier par un calcul, puis proposer une temporisation arrondie au dixième de seconde supérieur. \pts{1 pt}
  \rep{3.0}
\end{enumerate}


\end{document}
